\chapter{Orbitales de fragmentos y moléculas poliatómicas}\label{ch:b2:fragment-orbitals}

La estructura de Lewis del agua muestra dos enlaces O--H equivalentes y dos
pares no enlazantes equivalentes. Sin embargo, cuando el agua se ioniza con luz ultravioleta, los
electrones salen con cuatro energías distintas, no dos; el metano, con cuatro
enlaces C--H equivalentes, muestra dos. Los electrones de una molécula no se sitúan en enlaces
entre pares de átomos: ocupan orbitales extendidos por toda la molécula.
Este capítulo construye esos orbitales a partir de los \omterm{def:b2:fragment-orbitals:fragment}{orbitales de fragmentos}, explica
por qué el agua es angular y el eteno plano, y por qué un carbocatión prefiere tener
grupos alquilo a su alrededor.

\begin{recall}
\Cref{ch:b2:diatomic-mos}: el método \omterm{def:b2:diatomic-mos:lcao}{CLOA}, solapamiento y simetría, dos niveles que interaccionan,
\omterm{def:b2:diatomic-mos:bonding}{orbitales enlazantes} y antienlazantes. El volumen del primer año: geometrías RPECV, la hibridación
como descripción, la resonancia, los carbocationes y su orden de estabilidad (allí se anunció la
interacción de los enlaces C--H con el orbital vacío).
\end{recall}

\section{El método de los fragmentos}

\begin{definition}[Fragmentos]\label{def:b2:fragment-orbitals:fragment}
Un \emph{fragmento}\index{fragmento} es un grupo de átomos de una molécula considerado
por separado (un átomo, una pareja de hidrógenos, un grupo \ce{CH2}), y un \emph{orbital de
fragmento}\index{orbital de fragmento}, uno de sus orbitales. Los orbitales de la
molécula se construyen haciendo interaccionar los orbitales de dos fragmentos, dos a
dos, como los orbitales atómicos de una molécula diatómica.
\end{definition}

\begin{definition}[Combinación adaptada a la simetría]\label{def:b2:fragment-orbitals:symmetry-adapted}
Una \emph{combinación adaptada a la simetría}\index{combinacion adaptada a la simetria@combinación adaptada a la simetría} de
orbitales atómicos equivalentes es una combinación que cada operación de simetría de la
molécula (una reflexión en un plano especular, una rotación alrededor de un eje) o bien
deja invariante o bien convierte en su opuesta, o, para los conjuntos degenerados, transforma
en combinaciones del mismo conjunto.
\end{definition}

\begin{proposition}[Dos hidrógenos]\label{prop:b2:fragment-orbitals:h2-pair}
Para dos hidrógenos equivalentes con orbitales 1s $h_1$ y $h_2$, las
\omterm{def:b2:fragment-orbitals:symmetry-adapted}{combinaciones adaptadas a la simetría} son $h_1 + h_2$ y $h_1 - h_2$: la primera no
cambia, y la segunda cambia de signo, con cualquier operación que intercambie los dos
átomos.
\end{proposition}

\begin{proof}
Una operación que intercambia los átomos convierte $h_1$ en $h_2$ y $h_2$ en
$h_1$: $h_1 + h_2$ no cambia y $h_1 - h_2$ se convierte en $h_2 - h_1 = -(h_1 - h_2)$.
\end{proof}

\begin{proposition}[Tres y cuatro hidrógenos]\label{prop:b2:fragment-orbitals:h4}
Para cuatro hidrógenos en los vértices de un tetraedro, las combinaciones son una
combinación totalmente simétrica $h_1 + h_2 + h_3 + h_4$ y un conjunto de tres combinaciones
degeneradas con un plano nodal cada una, con forma de orbitales p. Para tres
hidrógenos en triángulo (como en el amoníaco), una combinación simétrica y una
pareja degenerada.
\end{proposition}

\admitted

Las formas de las combinaciones degeneradas, y sus nombres (a, e, t), proceden de la
teoría de grupos, que se desarrolla en el volumen del tercer año; aquí los nombres son solo etiquetas.

\begin{proposition}[Solo interacciona lo semejante]\label{prop:b2:fragment-orbitals:interaction-rules}
Un \omterm{def:b2:fragment-orbitals:fragment}{orbital de fragmento} solo interacciona con los orbitales del otro \omterm{def:b2:fragment-orbitals:fragment}{fragmento} que
se comportan del mismo modo frente a todas las operaciones de simetría de la molécula; la
intensidad de la interacción crece con el solapamiento y disminuye con la diferencia
de energía.
\end{proposition}

\begin{proof}
Si dos orbitales se comportan de modo distinto frente a alguna reflexión, uno es simétrico y
el otro antisimétrico respecto de ese plano, y su solapamiento y su
\omterm{def:b2:diatomic-mos:integrals}{integral de resonancia} se anulan (\cref{prop:b2:diatomic-mos:symmetry}). La dependencia del
solapamiento y de la diferencia de energía es la del \cref{thm:b2:diatomic-mos:heteronuclear}.
\end{proof}

\begin{method}[Un diagrama de fragmentos]\label{met:b2:fragment-orbitals:fragment-diagram}
\begin{enumerate}
  \item Divídase la molécula en dos \omterm{def:b2:fragment-orbitals:fragment}{fragmentos}, normalmente un átomo central y el conjunto
        de los átomos que lo rodean.
  \item Constrúyanse las \omterm{def:b2:fragment-orbitals:symmetry-adapted}{combinaciones adaptadas a la simetría} de los orbitales externos.
  \item Clasifíquense los orbitales de ambos \omterm{def:b2:fragment-orbitals:fragment}{fragmentos} según su comportamiento frente a los
        planos especulares y los ejes de la molécula.
  \item Hágase interaccionar cada orbital con los del mismo comportamiento, con más fuerza
        cuanto más próximos estén en energía; los restantes quedan no enlazantes.
  \item Llénese con los electrones de valencia y léase: nivel ocupado más alto,
        carácter enlazante, preferencias de forma.
\end{enumerate}
\end{method}

\section{El agua}

Sitúese el agua en el plano $yz$, con el eje $z$ como bisectriz del ángulo H--O--H. Dos
planos especulares contienen el eje $z$: el plano molecular $yz$ y el plano $xz$,
perpendicular a él. Respecto del plano $xz$, que intercambia los dos
hidrógenos, $h_1 + h_2$ es simétrica y $h_1 - h_2$ antisimétrica. En el oxígeno, los orbitales 2s y
$2\mathrm p_z$ son simétricos respecto de ambos planos; $2\mathrm p_y$ (en el
plano molecular, perpendicular a $z$) es antisimétrico respecto de $xz$;
$2\mathrm p_x$ es antisimétrico respecto del plano molecular, cosa que no es ninguna
combinación de los hidrógenos.

Así, $h_1 + h_2$ se mezcla con 2s y $2\mathrm p_z$, $h_1 - h_2$ con $2\mathrm p_y$, y
$2\mathrm p_x$ queda no enlazante. Resultan seis orbitales, cuatro de ellos llenos con los ocho
electrones de valencia.

\begin{omfigure}
\begin{tikzpicture}[font=\small]
  % O fragment (left)
  \pic (os) at (0,-2.4) {omlevel={ud}{0.7}}; \node[left] at (-0.45,-2.4) {2s};
  \pic (op1) at (-0.45,0) {omlevel={ud}{0.4}}; \pic (op2) at (0,0) {omlevel={u}{0.4}}; \pic (op3) at (0.45,0) {omlevel={u}{0.4}};
  \node[left] at (-0.7,0) {2p};
  \node at (0,-3.2) {O};
  % H2 fragment (right)
  \pic (hp) at (6,0.9) {omlevel={u}{0.7}}; \node[right] at (6.4,0.9) {$h_1 + h_2$};
  \pic (hm) at (6,1.8) {omlevel={u}{0.7}}; \node[right] at (6.4,1.8) {$h_1 - h_2$};
  \node[anchor=north west] at (6.4,0.6) {fragmento \ce{H2}};
  % MOs
  \pic (m1) at (3,-2.9) {omlevel={ud}{0.7}};  \node[omsymlabel, below=0.17cm, inner sep=0.5pt] at (m1-c) {$2a_1$};
  \pic (m2) at (3,-1.55) {omlevel={ud}{0.7}};  \node[omsymlabel, below=0.17cm, inner sep=0.5pt] at (m2-c) {$1b_2$};
  \pic (m3) at (3,-0.78) {omlevel={ud}{0.7}};  \node[omsymlabel, below=0.17cm, inner sep=0.5pt] at (m3-c) {$3a_1$};
  \pic (m4) at (3,0) {omlevel={ud}{0.7}};  \node[omsymlabel, below=0.17cm, inner sep=0.5pt] at (m4-c) {$1b_1$};
  \pic (m5) at (3,2.4) {omlevel={empty}{0.7}}; \node[omsymlabel, below=0.17cm, inner sep=0.5pt] at (m5-c) {$4a_1$};
  \pic (m6) at (3,3.1) {omlevel={empty}{0.7}}; \node[omsymlabel, below=0.17cm, inner sep=0.5pt] at (m6-c) {$2b_2$};
  \draw[omcorr, densely dashed] (os-r) -- (m1-l) (op3-r) -- (m2-l) (op3-r) -- (m3-l) (op3-r) -- (m4-l) (op3-r) -- (m5-l) (op3-r) -- (m6-l);
  \draw[omcorr, densely dashed] (hp-l) -- (m1-r) (hm-l) -- (m2-r) (hp-l) -- (m3-r) (hp-l) -- (m5-r) (hm-l) -- (m6-r);
\end{tikzpicture}

{\small Diagrama de \omterm{def:b2:fragment-orbitals:fragment}{fragmentos} del agua, esquema. La combinación simétrica de los
hidrógenos se mezcla con los orbitales 2s y $2\mathrm p_z$ del oxígeno (orbitales $a_1$),
y la antisimétrica, con $2\mathrm p_y$ ($b_2$); $2\mathrm p_x$, perpendicular al
plano molecular, queda no enlazante ($1b_1$, al nivel del 2p del oxígeno). Cuatro niveles de valencia
llenos, cuatro bandas fotoelectrónicas; la más alta, $1b_1$, se ioniza a
\qty{12.62}{eV}. Las etiquetas
$a_1$, $b_1$, $b_2$ son nombres (de la teoría de grupos, en el volumen del tercer año).}
% ledger: ie:H2O
\end{omfigure}

\begin{proposition}[Por qué el agua es angular]\label{prop:b2:fragment-orbitals:bent-water}
En un H--O--H lineal, dos orbitales p del oxígeno perpendiculares al eje son
no enlazantes. Al doblarse la molécula, uno de ellos, el del plano molecular, puede solaparse
con $h_1 + h_2$ y mezclarse con el orbital 2s: ese nivel lleno baja, más de lo que
suben los demás niveles llenos. Con ocho electrones de valencia la molécula angular es
la más estable.
\end{proposition}

\begin{proof}[Argumento]
En la molécula lineal ($z$ a lo largo de H--O--H), $h_1 + h_2$ tiene solapamiento nulo con
$2\mathrm p_x$ y $2\mathrm p_y$ (cada uno es antisimétrico respecto de un plano
que contiene el eje, y la combinación es simétrica); son dos pares degenerados
no enlazantes. Al doblarse en el plano $yz$, $2\mathrm p_y$ apunta en parte hacia
los hidrógenos: su solapamiento con $h_1 + h_2$ crece desde cero, y el nivel lleno
que lleva se estabiliza (el nivel $3a_1$ de la molécula angular). El nivel enlazante
$1b_2$ pierde algo de solapamiento y sube ligeramente. El balance es un descenso,
máximo cuando el nivel que baja está lleno, como ocurre con ocho electrones. Una
molécula con solo cuatro electrones de valencia, como \ce{BeH2}, tiene ese nivel
vacío y es lineal.
\end{proof}

\section{Metano, amoníaco y espectros fotoelectrónicos}

En el metano los cuatro orbitales de los hidrógenos dan la combinación simétrica, que
corresponde al orbital 2s del carbono, y un conjunto degenerado de tres, que corresponde a
sus tres orbitales 2p. Se llenan cuatro \omterm{def:b2:diatomic-mos:bonding}{orbitales enlazantes}: uno bajo ($a_1$) y tres
degenerados a mayor energía ($t_2$). El amoníaco, con tres hidrógenos, conserva un
orbital lleno sobre el nitrógeno que apunta en dirección opuesta a los hidrógenos, su par no enlazante, como
nivel ocupado más alto.

\begin{omfigure}
\begin{tikzpicture}[font=\small]
  \pic (cs) at (0,-1.6) {omlevel={ud}{0.7}}; \node[left] at (-0.45,-1.6) {2s};
  \pic (cp1) at (-0.45,0) {omlevel={u}{0.4}}; \pic (cp2) at (0,0) {omlevel={u}{0.4}}; \pic (cp3) at (0.45,0) {omlevel={empty}{0.4}};
  \node[left] at (-0.7,0) {2p};
  \node at (0,-2.5) {C};
  \pic (ha) at (6,-0.8) {omlevel={ud}{0.7}}; \node[right] at (6.4,-0.8) {$a_1$};
  \pic (ht1) at (5.55,0.4) {omlevel={u}{0.4}}; \pic (ht2) at (6,0.4) {omlevel={u}{0.4}}; \pic (ht3) at (6.45,0.4) {omlevel={empty}{0.4}};
  \node[right] at (6.7,0.4) {$t_2$};
  \node at (6,-1.6) {fragmento \ce{H4}};
  \pic (m1) at (3,-2.4) {omlevel={ud}{0.7}}; \node[omsymlabel, below=0.17cm, inner sep=0.5pt] at (m1-c) {$1a_1$};
  \pic (m2a) at (2.55,-0.9) {omlevel={ud}{0.4}}; \pic (m2b) at (3,-0.9) {omlevel={ud}{0.4}}; \pic (m2c) at (3.45,-0.9) {omlevel={ud}{0.4}};
  \node[omsymlabel, below=0.17cm, inner sep=0.5pt] at (m2b-c) {$1t_2$};
  \pic (m3) at (3,1.8) {omlevel={empty}{0.7}}; \node[omsymlabel, below=0.17cm, inner sep=0.5pt] at (m3-c) {$2a_1$};
  \pic (m4a) at (2.55,2.5) {omlevel={empty}{0.4}}; \pic (m4b) at (3,2.5) {omlevel={empty}{0.4}}; \pic (m4c) at (3.45,2.5) {omlevel={empty}{0.4}};
  \node[omsymlabel, below=0.17cm, inner sep=0.5pt] at (m4b-c) {$2t_2$};
  \draw[omcorr, densely dashed] (cs-r) -- (m1-l) (cp3-r) -- (m2a-l) (cs-r) -- (m3-l) (cp3-r) -- (m4a-l);
  \draw[omcorr, densely dashed] (ha-l) -- (m1-r) (ht1-l) -- (m2c-r) (ha-l) -- (m3-r) (ht1-l) -- (m4c-r);
\end{tikzpicture}

{\small Diagrama de \omterm{def:b2:fragment-orbitals:fragment}{fragmentos} del metano, esquema. La combinación simétrica de los
cuatro hidrógenos corresponde al orbital 2s del carbono, y las tres degeneradas, a sus orbitales
2p: dos niveles llenos, y por tanto dos bandas fotoelectrónicas, la superior
(que se ioniza primero, a \qty{12.61}{eV}) con seis electrones.}
% ledger: ie:CH4
\end{omfigure}

\begin{definition}[Espectroscopia fotoelectrónica]\label{def:b2:fragment-orbitals:photoelectron}
La \emph{espectroscopia fotoelectrónica}\index{espectroscopia fotoelectronica@espectroscopia fotoelectrónica} mide las
energías cinéticas de los electrones arrancados de las moléculas por una radiación
monocromática de energía conocida; cada banda del espectro corresponde a la extracción
de un electrón de un nivel ocupado, y su posición da la energía de ionización
de ese nivel.
\end{definition}

\begin{proposition}[Aproximación de Koopmans]\label{prop:b2:fragment-orbitals:koopmans}
La energía necesaria para arrancar un electrón de un orbital ocupado es
aproximadamente igual a menos la energía de ese orbital, $I \approx -\varepsilon$.
\end{proposition}

\admitted

La aproximación desprecia la relajación de los demás electrones tras la
ionización; se justifica en el volumen del tercer año. Convierte un espectro
fotoelectrónico en un diagrama de orbitales: el agua muestra cuatro bandas de valencia, y el metano dos,
de acuerdo con los diagramas de \omterm{def:b2:fragment-orbitals:fragment}{fragmentos}, y en contra de la imagen de dos pares no enlazantes
equivalentes y dos enlaces equivalentes en el agua, que describe la misma densidad electrónica
total de otra manera pero no da las energías de los
electrones arrancados.

\section{El eteno}

El eteno puede construirse a partir de dos \omterm{def:b2:fragment-orbitals:fragment}{fragmentos} \ce{CH2}. Cada \ce{CH2} tiene, además de los
orbitales de sus dos enlaces C--H, un orbital en el plano que apunta en dirección opuesta a los hidrógenos
(y forma el enlace $\sigma$ C--C con su pareja) y un orbital p
perpendicular al plano. Los dos orbitales p se combinan en un $\pi$ lleno y un
$\pi^*$ vacío.

\begin{omfigure}
\begin{tikzpicture}[font=\small]
  \begin{scope}
    \pic at (0,0) {omp={0.85}{90}}; \pic at (1.2,0) {omp={0.85}{90}};
    \draw[thick] (0,0) -- (1.2,0);
    \draw[thick] (0,0) -- (-0.5,-0.45) (0,0) -- (-0.55,0.35) (1.2,0) -- (1.7,-0.45) (1.2,0) -- (1.75,0.35);
    \node at (0.6,-1.1) {$\pi$ (lleno)};
  \end{scope}
  \begin{scope}[xshift=3.6cm]
    \pic at (0,0) {omp={0.85}{90}}; \pic at (1.2,0) {omp={-0.85}{90}};
    \draw[thick] (0,0) -- (1.2,0);
    \draw[thick] (0,0) -- (-0.5,-0.45) (0,0) -- (-0.55,0.35) (1.2,0) -- (1.7,-0.45) (1.2,0) -- (1.75,0.35);
    \node at (0.6,-1.1) {$\pi^*$ (vacío)};
  \end{scope}
  \begin{scope}[xshift=7.4cm]
    \pic at (0,0) {omp={0.85}{90}}; \pic at (1.2,0) {omp={0.85}{0}};
    \draw[thick] (0,0) -- (1.2,0);
    \node at (0.6,-1.1) {girado $90^\circ$};
    \node[font=\scriptsize, align=center] at (0.6,1.3) {solapamiento nulo};
  \end{scope}
\end{tikzpicture}

{\small El sistema $\pi$ del eteno, esquema (la molécula vista de canto, con los
hidrógenos en el plano). Los dos orbitales p perpendiculares al plano se combinan
en fase ($\pi$, lleno) y en oposición de fase ($\pi^*$, vacío). Girar un \ce{CH2}
$90^\circ$ hace perpendiculares los dos orbitales p: su solapamiento, y el enlace $\pi$,
se anulan.}
\end{omfigure}

El enlace $\pi$ mantiene plano el eteno: girar un extremo aparta los orbitales p
uno del otro, su solapamiento disminuye como el coseno del ángulo de giro y, a $90^\circ$,
el enlace $\pi$ ha desaparecido. Para girar la molécula hay que aportar la energía de un enlace $\pi$,
y por eso los isómeros cis y trans de los alquenos no se interconvierten a
temperatura ambiente, mientras que la rotación alrededor de un enlace sencillo es libre.

\section{Hiperconjugación}

\begin{definition}[Hiperconjugación]\label{def:b2:fragment-orbitals:hyperconjugation}
La \emph{hiperconjugación}\index{hiperconjugacion@hiperconjugación} es la interacción de un orbital
$\sigma$ lleno de un enlace C--H o C--C con un orbital p o $\pi^*$ vecino, vacío o parcialmente
lleno, paralelo a él.
\end{definition}

\begin{proposition}[Los grupos alquilo estabilizan los carbocationes]\label{prop:b2:fragment-orbitals:hyperconjugation}
Un carbocatión está estabilizado por cada enlace C--H o C--C de un carbono vecino que
puede alinearse con su orbital p vacío: cuantos más enlaces de este tipo, más estable es el
catión, de ahí el orden terciario > secundario > primario > metilo.
\end{proposition}

\begin{proof}[Argumento]
El orbital p vacío del carbono catiónico y un orbital $\sigma_{\text{C--H}}$ lleno
contiguo forman un sistema de dos niveles con una gran diferencia de energía $\Delta$ y un acoplamiento pequeño
$\beta$ (no nulo solo si el enlace es aproximadamente paralelo al orbital p). Por el
\cref{thm:b2:diatomic-mos:heteronuclear}, el nivel $\sigma$ lleno baja en unos
$\beta^2/\Delta$, y los dos electrones ganan el doble; el nivel vacío sube, sin
coste. Cada enlace convenientemente alineado aporta su propia estabilización; un grupo metilo
siempre tiene un enlace C--H casi alineado, sea cual sea su rotación.
\end{proof}

\begin{omfigure}
\begin{tikzpicture}[font=\small]
  \pic at (0,0) {omp={0.9}{90}};
  \node[circle, fill=black, inner sep=1.2pt] at (0,0) {};
  \node[font=\scriptsize, anchor=north west] at (0.08,-0.05) {C$^+$};
  \draw[thick] (0,0) -- (1.4,0);
  \node[circle, fill=black, inner sep=1.2pt] at (1.4,0) {};
  \draw[thick] (1.4,0) -- (1.75,0.85);
  \pic at (1.58,0.42) {omlobe={0.75}{68}};
  \pic at (1.58,0.42) {omlobe={0.45}{248}};
  \node[font=\scriptsize] at (1.95,1.0) {H};
  \draw[thick] (1.4,0) -- (2.0,-0.45) (1.4,0) -- (1.1,-0.6);
  \draw[<->, omProp, thick] (0.35,0.8) .. controls (0.8,1.2) and (1.2,1.2) .. (1.5,0.95);
  \node[font=\scriptsize, align=left, anchor=west] at (2.4,0.4) {lóbulo $\sigma_{\text{C--H}}$ lleno paralelo\\ al orbital p vacío};
  % level scheme
  \begin{scope}[xshift=7.2cm, yshift=-0.4cm]
    \pic (p) at (0,1.4) {omlevel={empty}{0.7}}; \node[left] at (-0.4,1.4) {p vacío};
    \pic (s) at (2.2,-0.8) {omlevel={ud}{0.7}}; \node[right] at (2.6,-0.8) {$\sigma_{\text{C--H}}$};
    \pic (lo) at (1.1,-1.2) {omlevel={ud}{0.7}};
    \pic (hi) at (1.1,1.7) {omlevel={empty}{0.7}};
    \draw[omcorr, densely dashed] (p-r) -- (lo-l) (p-r) -- (hi-l) (s-l) -- (lo-r) (s-l) -- (hi-r);
    \node[font=\scriptsize, anchor=north] at (1.1,-1.35) {estabilizado en $\approx \beta^2/\Delta$};
  \end{scope}
\end{tikzpicture}

{\small \omterm{def:b2:fragment-orbitals:hyperconjugation}{Hiperconjugación} en un carbocatión, esquema. Un \omterm{def:b2:diatomic-mos:bonding}{orbital enlazante} C--H lleno
del carbono vecino, paralelo al orbital p vacío, le cede parte de su
densidad electrónica; el nivel lleno se estabiliza como en cualquier interacción de
dos niveles de energías distintas.}
\end{omfigure}

\section{Ejercicios}

\begin{exercise}[$\star$]\label{exo:b2:fragment-orbitals:1}
Escríbanse las dos \omterm{def:b2:fragment-orbitals:symmetry-adapted}{combinaciones adaptadas a la simetría} de los orbitales 1s de los hidrógenos del
agua y dígase cuál de ellas cambia de signo en el plano bisector del ángulo H--O--H
perpendicular a la molécula.
\end{exercise}

\begin{exercise}[$\star$]\label{exo:b2:fragment-orbitals:2}
¿Cuántos \omterm{def:b2:diatomic-mos:lcao}{orbitales moleculares} de valencia tiene el agua, cuántos electrones de valencia
y cuántos orbitales de valencia llenos?
\end{exercise}

\begin{exercise}[$\star$]\label{exo:b2:fragment-orbitals:3}
¿Por qué el espectro fotoelectrónico del metano muestra dos bandas de valencia, una de ellas
tres veces más intensa que la otra?
\end{exercise}

\begin{exercise}[$\star$]\label{exo:b2:fragment-orbitals:4}
¿Por qué los seis átomos del eteno están en un mismo plano?
\end{exercise}

\begin{exercise}[$\star\star$]\label{exo:b2:fragment-orbitals:5}
La primera energía de ionización del amoníaco es de \qty{10.07}{eV}, y la del metano, de
\qty{12.61}{eV}. ¿Qué orbital pierde el electrón en cada caso, y por qué es el amoníaco
el más fácil de ionizar?
\end{exercise}

\begin{exercise}[$\star\star$]\label{exo:b2:fragment-orbitals:6}
Compárense los niveles ocupados del agua lineal y del agua angular, y explíquese por qué una molécula
como \ce{BeH2}, con cuatro electrones de valencia, es lineal.
\end{exercise}

\begin{exercise}[$\star\star$]\label{exo:b2:fragment-orbitals:7}
Explíquese, con el resultado de los dos niveles, el orden de estabilidad de los carbocationes
\ce{CH3+}, \ce{CH3CH2+}, \ce{(CH3)2CH+}, \ce{(CH3)3C+}.
\end{exercise}

\begin{exercise}[$\star\star$]\label{exo:b2:fragment-orbitals:8}
El solapamiento de dos orbitales p girados un ángulo $\theta$ alrededor del enlace es
proporcional a $\cos\theta$. Con una \omterm{def:b2:diatomic-mos:integrals}{integral de resonancia} proporcional al
solapamiento, ¿cómo depende la estabilización $\pi$ de $\theta$? ¿A qué ángulo
se reduce a la mitad?
\end{exercise}

\begin{exercise}[$\star\star$]\label{exo:b2:fragment-orbitals:9}
El borano \ce{BH3} es plano, y el amoníaco, piramidal. Con los \omterm{def:b2:fragment-orbitals:fragment}{orbitales de fragmento} de un
átomo central y de tres hidrógenos, explíquese la diferencia por los electrones que tiene cada uno.
\end{exercise}

\begin{exercise}[$\star\star\star$]\label{exo:b2:fragment-orbitals:10}
El ion \ce{H3+} es un triángulo equilátero. Con $\alpha$ y $\beta$ para los orbitales 1s
y despreciando el solapamiento, hállense sus tres \omterm{def:b2:atomic-orbitals:orbital-energy}{energías orbitales} (la combinación simétrica
tiene energía $\alpha + 2\beta$) y explíquese por qué bastan dos electrones para
mantenerlo ligado.
\end{exercise}

\begin{exercise}[$\star\star\star$]\label{exo:b2:fragment-orbitals:11}
El sistema alilo \ce{CH2=CH-CH2} tiene tres orbitales p paralelos. Adivínense las formas de
sus tres orbitales $\pi$ (signos en cada carbono) a partir del número de nodos, y cuál
es el nivel ocupado más alto del anión alilo.
\end{exercise}

\begin{exercise}[$\star\star\star$]\label{exo:b2:fragment-orbitals:12}
Constrúyase el sistema $\pi$ del dióxido de carbono a partir de los orbitales 2p perpendiculares al
eje O--C--O: ¿cuántos orbitales $\pi$ hay, cuáles son enlazantes, no enlazantes y
antienlazantes, y cuántos electrones contienen?
\end{exercise}

\section{Problema: la forma del agua}

\begin{problem}[{Problema de fin de semana --- los \omterm{def:b2:fragment-orbitals:fragment}{fragmentos} oxígeno y \ce{H2}, los orbitales de
un H--O--H lineal, los niveles que se desplazan al doblarse y el espectro
fotoelectrónico con la aproximación de Koopmans}]
\label{pb:b2:fragment-orbitals:1}
Datos: geometría del agua, $r(\ce{O-H}) = \qty{95.8}{pm}$, ángulo
$104.48^\circ$; primeras energías de ionización: \ce{H2O} \qty{12.62}{eV}, \ce{O}
\qty{13.62}{eV}. El agua está en el plano $yz$, con $z$ a lo largo de la bisectriz.
% ledger: geo:H2O.rOH, geo:H2O.aHOH, ie:H2O, ie:O

\medskip
\noindent\textbf{Parte I --- \omterm{def:b2:fragment-orbitals:fragment}{Fragmentos}.}
\begin{enumerate}
  \item ¿Cuántos electrones de valencia tiene el agua, y de qué átomos proceden?
  \item Escríbanse las dos \omterm{def:b2:fragment-orbitals:symmetry-adapted}{combinaciones adaptadas a la simetría} de los orbitales 1s de los hidrógenos.
  \item ¿Qué orbitales del oxígeno son simétricos respecto de ambos planos especulares?
  \item ¿Qué orbital del oxígeno es antisimétrico solo respecto del plano $xz$?
  \item ¿Qué orbital del oxígeno no tiene pareja entre las combinaciones de los hidrógenos?
  \item ¿Cuántos \omterm{def:b2:diatomic-mos:lcao}{orbitales moleculares} resultan, y cuántos están llenos?
\end{enumerate}

\medskip
\noindent\textbf{Parte II --- Un H--O--H lineal.}
\begin{enumerate}[resume]
  \item En una molécula lineal a lo largo de $z$, ¿qué combinación interacciona con el
        $2\mathrm p_z$ del oxígeno?
  \item ¿Cuál interacciona con el 2s del oxígeno?
  \item ¿Qué orbitales del oxígeno quedan no enlazantes, y con qué degeneración?
  \item Llénense los niveles con los ocho electrones de valencia.
  \item ¿Sería \omterm{def:b2:diatomic-mos:magnetism}{paramagnética} el agua lineal?
  \item ¿Dónde estaría su nivel ocupado más alto, comparado con un orbital 2p del
        oxígeno?
\end{enumerate}

\medskip
\noindent\textbf{Parte III --- La flexión.}
\begin{enumerate}[resume]
  \item La molécula lineal situada a lo largo de $z$ se dobla en el plano $yz$. ¿Qué
        \omterm{def:b2:diatomic-mos:bonding}{orbital no enlazante} empieza a solaparse con $h_1 + h_2$?
  \item ¿Qué le ocurre al nivel lleno que lleva?
  \item ¿Qué nivel enlazante sube ligeramente, y por qué?
  \item Conclúyase: ¿por qué es angular el agua?
  \item Compárese el ángulo medido con $90^\circ$ (enlace p puro) y con $109.5^\circ$.
  \item ¿Qué \omterm{def:b2:diatomic-mos:bonding}{orbital no enlazante} no se ve afectado por la flexión?
\end{enumerate}

\medskip
\noindent\textbf{Parte IV --- El espectro fotoelectrónico.}
\begin{enumerate}[resume]
  \item ¿Cuántas bandas de valencia muestra el espectro?
  \item Según la aproximación de Koopmans, ¿de qué orbital procede la primera banda?
  \item ¿Por qué esa banda es estrecha, mientras que las bandas de los \omterm{def:b2:diatomic-mos:bonding}{orbitales enlazantes} son anchas?
  \item Compárese la primera energía de ionización del agua con la del átomo de
        oxígeno, y explíquese la diferencia.
  \item ¿Qué queda de los ``dos pares no enlazantes equivalentes'' de la estructura de Lewis?
  \item Indíquense el número de bandas de ionización de valencia del agua y la energía de
        la más baja.
\end{enumerate}
\end{problem}
